%%
% The BIThesis Template for Undergraduate Thesis
%
% 北京理工大学毕业设计（论文） —— 使用 XeLaTeX 编译
%
% Copyright 2021-2023 BITNP
%
% This work may be distributed and/or modified under the
% conditions of the LaTeX Project Public License, either version 1.3
% of this license or (at your option) any later version.
% The latest version of this license is in
%   http://www.latex-project.org/lppl.txt
% and version 1.3 or later is part of all distributions of LaTeX
% version 2005/12/01 or later.
%
% This work has the LPPL maintenance status `maintained'.
%
% The Current Maintainer of this work is Feng Kaiyu.
%
% Compile with: xelatex -> biber -> xelatex -> xelatex

% !TeX program = xelatex
% !BIB program = biber


% 开启盲审格式 blindPeerReview=true (如：[type=bachelor,blindPeerReview=true])

\documentclass[type=bachelor]{thesis}

% 此处仅列出常用的配置。全部配置用法请见「bithesis.pdf」手册。
\BITSetup{
  cover = {
    % 封面需要「北京理工大学」字样图片，如无必要请勿修改该项。
    headerImage = images/header.png,
    % 封面标题需要“华文细黑”，如无必要请勿修改该项。
    xiheiFont = STXIHEI.TTF,
    % 可用以下参数自定义封面日期。
    % date = 1955年11月,
  },
  info = {
    % 想要删除某项封面信息，直接删除该项即可。
    % 想要让某项封面信息留空（但是保留下划线），请传入空白符组成的字符串，如"{~}"。
    % 如需要换行，则用 “\\” 符号分割。
    title = 北京理工大学本科生毕业设计（论文）题目,
    titleEn = {The Subject of Undergraduate Graduation Project (Thesis) of Beijing Institute of Technology},
    school = 计算机学院,
    major = 计算机科学与技术,
    class = 0561xxxx,
    author = 惠计算,
    studentId = 11xxxxxxxx,
    supervisor = 张哈希,
    keywords = {北京理工大学；本科生；毕业设计（论文）——请在“main.tex”开头设置},
    keywordsEn = {BIT; Undergraduate; Graduation Project (Thesis)},
    % 如果你的毕设为校外毕设，请将下面这一行语句解除注释（删除第一个百分号字符）并填写你的校外毕设导师名字
    % externalSupervisor = 左偏树,
  },
  style = {
    % 开启 Windows 平台下的中易宋体伪粗体。
    % windowsSimSunFakeBold = true,
    %
    % 开启该选项后，将用 Times New Roman 的开源字体 TeX Gyre Termes 作为正文字体。
    % 这个选项适用于以下情况：
    % 1. 不想在系统中安装 Times New Roman。
    % 2. 在 Linux/macOS 下遇到 `\textsc` 无法正常显示的问题。
    % betterTimesNewRoman = true,
  },
  misc = {
    % 微调表格行间距
    tabularRowSeparation = 1.25,
  },
}

% 使用 listings 宏包进行代码块使用，并使用了预定义的样式，
% 你也可以选用自己的喜欢的其他宏包，如 minted；
% 然而由于 minted 依赖 Python 的 Pygments 库作为外部依赖，因此出于模板的简洁程度考虑，我们没有提供 minted 进行代码块书写的示例。
\usepackage{listings}


% 大部分关于参考文献样式的修改，都可以通过此处的选项进行配置。
% 详情请搜索「biblatex-gb7714-2015 文档」进行阅读。
\usepackage[
  backend=biber,
  style=gb7714-2015,
  gbalign=gb7714-2015,
  gbnamefmt=lowercase,
  gbpub=false,
  doi=false,
  url=false,
  eprint=false,
  isbn=false,
]{biblatex}

% 参考文献引用文件位于 misc/ref.bib
\addbibresource{misc/ref.bib}

% 如果要按照计算机学院的要求，带有“北京理工大学”水印，请使用此 issue 提供的方法：
% https://github.com/BITNP/BIThesis/issues/350#issuecomment-1565974141

% 文档开始
\begin{document}

% 标题页面：如无特殊需要，本部分无需改动
\MakeCover

% 原创性声明：盲审结束前无需改动
% 后续添加签名时，可考虑先用 Word 制作再覆盖 misc/1_originality.pdf。
% 不过如无特殊需要，本部分的代码仍无需改动。
\begin{blindPeerReview}
  \includepdf{misc/1_originality.pdf}\newpage
\end{blindPeerReview}
% %%
% The BIThesis Template for Bachelor Graduation Thesis
%
% 北京理工大学毕业设计（论文）原创性声明模板 —— 使用 XeLaTeX 编译
%
% Copyright 2020-2023 BITNP
%
% This work may be distributed and/or modified under the
% conditions of the LaTeX Project Public License, either version 1.3
% of this license or (at your option) any later version.
% The latest version of this license is in
%   http://www.latex-project.org/lppl.txt
% and version 1.3 or later is part of all distributions of LaTeX
% version 2005/12/01 or later.
%
% This work has the LPPL maintenance status `maintained'.
%
% The Current Maintainer of this work is Feng Kaiyu.
%
% Compile with: xelatex -> biber -> xelatex -> xelatex
%
% 如无特殊需要，本页面无需更改

\MakeOriginality


% 前置内容定义
\frontmatter
% 摘要：在相应 TeX 文件撰写
%%
% The BIThesis Template for Bachelor Graduation Thesis
%
% 北京理工大学毕业设计（论文）中英文摘要 —— 使用 XeLaTeX 编译
%
% Copyright 2020-2023 BITNP
%
% This work may be distributed and/or modified under the
% conditions of the LaTeX Project Public License, either version 1.3
% of this license or (at your option) any later version.
% The latest version of this license is in
%   http://www.latex-project.org/lppl.txt
% and version 1.3 or later is part of all distributions of LaTeX
% version 2005/12/01 or later.
%
% This work has the LPPL maintenance status `maintained'.
%
% The Current Maintainer of this work is Feng Kaiyu.

% 中英文摘要章节
\begin{abstract}
% 中文摘要正文从这里开始
本文……。

\textcolor{blue}{摘要正文选用模板中的样式所定义的“正文”，每段落首行缩进 2 个字符；或者手动设置成每段落首行缩进 2 个汉字，字体：宋体，字号：小四，行距：固定值 22 磅，间距：段前、段后均为 0 行。阅后删除此段。}

\textcolor{blue}{摘要是一篇具有独立性和完整性的短文，应概括而扼要地反映出本论文的主要内容。包括研究目的、研究方法、研究结果和结论等，特别要突出研究结果和结论。中文摘要力求语言精炼准确，本科生毕业设计（论文）摘要建议 300-500 字。摘要中不可出现参考文献、图、表、化学结构式、非公知公用的符号和术语。英文摘要与中文摘要的内容应一致。阅后删除此段。}

\end{abstract}

% 一般情况下，超出该行宽度的最后一个单词会被排版引擎尝试（在适合的地方）插入连字符（hyphen）
% 换行。然而，当单词本身比较生僻的情况下，LaTeX 可能会保留完整单词，从而导致该行宽度
% 绕过限制。以“aaaaaabbb” 为例，你可以在以下两个方法中的一种来指定插入连字符的位置：
% 1. 在此处使用 \hyphenation{aaaaaa-bbb}。
% 2. 在文章中使用 aaaaaa\-bbb。
% 以上两种方法都能让引擎在 ab 交界处插入连字符，从而正常换行。

% 英文摘要章节
\begin{abstractEn}
% 英文摘要正文从这里开始
In order to study……

\textcolor{blue}{Abstract 正文设置成每段落首行缩进 2 字符，字体：Times New Roman，字号：小四，行距：固定值 22 磅，间距：段前、段后均为 0 行。阅后删除此段。}
\end{abstractEn}


\MakeTOC

% 正文开始
\mainmatter

% 在这里引用各章 TeX 文件，按需添加
%%
% The BIThesis Template for Bachelor Graduation Thesis
%
% 北京理工大学毕业设计（论文）第一章节 —— 使用 XeLaTeX 编译
%
% Copyright 2020-2023 BITNP
%
% This work may be distributed and/or modified under the
% conditions of the LaTeX Project Public License, either version 1.3
% of this license or (at your option) any later version.
% The latest version of this license is in
%   http://www.latex-project.org/lppl.txt
% and version 1.3 or later is part of all distributions of LaTeX
% version 2005/12/01 or later.
%
% This work has the LPPL maintenance status `maintained'.
%
% The Current Maintainer of this work is Feng Kaiyu.
%
% 第一章节

\chapter{一级题目}

\section{二级题目}
% 这里插入一个参考文献，仅作参考

\subsection{三级题目}

正文……\parencite{yuFeiJiZongTiDuoXueKeSheJiYouHuaDeXianZhuangYuFaZhanFangXiang2008}……\cite{Hajela2012Application}

\textcolor{blue}{正文部分：宋体、小四；正文行距：22磅；间距段前段后均为0行。阅后删除此段。}

\textcolor{blue}{图、表居中，图注标在图下方，表头标在表上方，宋体、五号、居中，1.25倍行距，间距段前段后均为0行，图表与上下文之间各空一行。阅后删除此段。}

\textcolor{blue}{\underline{\underline{图-示例：（阅后删除此段）}}}


\begin{figure}[htbp]
  \centering
  \includegraphics[]{images/bit_logo.png}
  \caption{标题序号}\label{标题序号} % label 用来在文中索引
\end{figure}

\textcolor{blue}{\underline{\underline{表-示例：（阅后删除此段）}}}
% 三线表
\begin{table}[htbp]
  \centering
  \caption{统计表}\label{统计表}
  \begin{tabular}{*{5}{>{\centering\arraybackslash}p{2cm}}} \toprule
    项目    & 产量    & 销量    & 产值   & 比重    \\ \midrule
    手机    & 1000  & 10000 & 500  & 50\%  \\
    计算机   & 5500  & 5000  & 220  & 22\%  \\
    笔记本电脑 & 1100  & 1000  & 280  & 28\%  \\ \midrule
    合计    & 17600 & 16000 & 1000 & 100\% \\ \bottomrule
    \end{tabular}
\end{table}

\textcolor{blue}{公式标注应于该公式所在行的最右侧。对于较长的公式只可在符号处（+、-、*、/、$\leqslant$ $\geqslant$ 等）转行。在文中引用公式时，在标号前加“式”，如式（1-2）。阅后删除此
段。}

\textcolor{blue}{公式-示例：（阅后删除此段）}
% 公式上下不要空行，置于同一个段落下即可，否则上下距离会出现高度不一致的问题
\begin{equation}
    LRI=1\ ∕\ \sqrt{1+{\left(\frac{{\mu }_{R}}{{\mu }_{s}}\right)}^{2}{\left(\frac{{\delta }_{R}}{{\delta }_{s}}\right)}^{2}}
\end{equation}

\subsubsection{生僻字}

% 一个可能无法正常显示的生僻字: 彧。下文注释中，介绍了如何通过自定义字体来显示生僻字。

% 定义一个提供了生僻字的字体，注意要确保你的系统存在该字体
% \setCJKfamilyfont{custom-font}{Noto Serif CJK SC}

% 使用自己定义的字体
% 使用提供了相应字形的字体：{\CJKfamily{custom-font} 彧}。
% \CJKfamily 会切换字体，影响之后所有内容。故另套“{}”来分组，限制其作用范围。


\section{字体效果表格}

% 列：Regular、Italic、Bold、Bold Italic
% 行：宋体、黑体、楷体、Serif、Sans Serif、Typewriter、Math

\begin{table}[htb]
    \centering
    \caption{字体效果表格}
    \begin{tabular}{@{}lllll@{}}
    \toprule
               & Regular & Bold & Italic & Bold Italic \\ \midrule
      宋体       & 宋体      & \colorbox{orange}{\textbf{宋体粗体}} & \textit{楷体}     &   \colorbox{gray}{\textbf{\textit{楷书粗斜体}}}  \\
      黑体         & {\heiti{}黑体}      & \textbf{\heiti{}黑体粗体} &   \textit{\heiti{}黑体斜体}     & \colorbox{gray}{\textit{\textbf{\heiti{}黑体粗斜体}}}   \\
      楷体         & {\kaishu{}楷书}      & \textbf{\kaishu{}楷书粗体} & \textit{\kaishu{}斜体楷体} &  \colorbox{gray}{\textbf{\textit{\kaishu{}楷书粗斜体}}}    \\
    Serif(Roman/Normal)      &    Regular    &  \textbf{Bold}  &    \textit{Italic}    &     \textbf{\textit{Bold Italic}}    \\
    Sans Serif &  \textsf{Regular}       &  \textbf{\textsf{Bold}}    &  \textit{\textsf{Bold}}   &    \textbf{\textit{\textsf{Bold}}}   \\
    % 有些字体缺少特定变体，LaTeX会抛出警告，同时尽量替换为相近变体。
    % 若编译结果能接受，可忽略之。
    % 例如此处Typewriter代表的Latin Modern Sans Typewriter（lmtt）字体没有bold extended italic（bx/it）变体，
    % 所以`\textbf{\textit{\texttt{…}}}`会被替换为bold slant（b/sl）变体，并引发如下警告。（其中TU指字体编码）
    % Font shape `TU/lmtt/bx/it' in size <…> not available. Font shape `TU/lmtt/b/sl' tried instead.
    Typewriter &  \texttt{Regular}       &  \textbf{\texttt{Bold}}    &  \textit{\texttt{Bold}}   &    \textbf{\textit{\texttt{Bold}}}   \\
    Math       &   $\mathnormal{Regular} \mathrm{Roman}$  & $\mathbf{Bold}$   &    $\mathit{Italic}$    &  $\mathbf{\mathit{Bold Italic}}$    \\ \bottomrule
    \end{tabular}
\end{table}

\begin{itemize}[nosep]
  \item \colorbox{orange}{宋体粗体}在 Windows 下会成为黑体。这是因为 Windows 的中易宋体没有粗体字重而进行的妥协。
    如果想要获得宋体粗体的样式，请在配置中开启伪粗体选项。
  \item \colorbox{gray}{粗斜体}的效果是因操作系统字体而定的，中文写作中不会使用这种字形，可以忽略。
\end{itemize}

\textit{有关公式与上下文间距的一些注意事项：请保证源码中的公式的环境（如}
\\ \verb|\begin{equation}|
  \textit{）与上一段落不要有空行。否则，公式和上文段落之间会有额外的空白。}

%%
% The BIThesis Template for Bachelor Graduation Thesis
%
% 北京理工大学毕业设计（论文）第二章节 —— 使用 XeLaTeX 编译
%
% Copyright 2020-2023 BITNP
%
% This work may be distributed and/or modified under the
% conditions of the LaTeX Project Public License, either version 1.3
% of this license or (at your option) any later version.
% The latest version of this license is in
%   http://www.latex-project.org/lppl.txt
% and version 1.3 or later is part of all distributions of LaTeX
% version 2005/12/01 or later.
%
% This work has the LPPL maintenance status `maintained'.
%
% The Current Maintainer of this work is Feng Kaiyu.
%%

\chapter{另一个章节}

\section{代码片段}

\begin{lstlisting}[language=Python, caption={Python Code}, label={lst:pythonfile}]
import numpy as np

def incmatrix(genl1,genl2):
    m = len(genl1)
    n = len(genl2)
    M = None #to become the incidence matrix
    VT = np.zeros((n*m,1), int)  #dummy variable

    #compute the bitwise xor matrix
    M1 = bitxormatrix(genl1)
    M2 = np.triu(bitxormatrix(genl2),1)

    for i in range(m-1):
        for j in range(i+1, m):
            [r,c] = np.where(M2 == M1[i,j])
            for k in range(len(r)):
                VT[(i)*n + r[k]] = 1;
                VT[(i)*n + c[k]] = 1;
                VT[(j)*n + r[k]] = 1;
                VT[(j)*n + c[k]] = 1;

                if M is None:
                    M = np.copy(VT)
                else:
                    M = np.concatenate((M, VT), 1)

                VT = np.zeros((n*m,1), int)

    return M
\end{lstlisting}


% 后置内容
\backmatter

% 结论：在相应 TeX 文件撰写
%%
% The BIThesis Template for Bachelor Graduation Thesis
%
% 北京理工大学毕业设计（论文）结论 —— 使用 XeLaTeX 编译
%
% Copyright 2020-2023 BITNP
%
% This work may be distributed and/or modified under the
% conditions of the LaTeX Project Public License, either version 1.3
% of this license or (at your option) any later version.
% The latest version of this license is in
%   http://www.latex-project.org/lppl.txt
% and version 1.3 or later is part of all distributions of LaTeX
% version 2005/12/01 or later.
%
% This work has the LPPL maintenance status `maintained'.
%
% The Current Maintainer of this work is Feng Kaiyu.
%
% Compile with: xelatex -> biber -> xelatex -> xelatex

\begin{conclusion}
  % 结论部分尽量不使用 \subsection 二级标题，只使用 \section 一级标题

  % 这里插入一个参考文献，仅作参考
  本文结论……\cite{张伯伟2002全唐五代诗格会考}。

  \textcolor{blue}{结论作为毕业设计（论文）正文的最后部分单独排写，但不加章号。结论是对整个论文主要结果的总结。在结论中应明确指出本研究的创新点，对其应用前景和社会、经济价值等加以预测和评价，并指出今后进一步在本研究方向进行研究工作的展望与设想。结论部分的撰写应简明扼要，突出创新性。阅后删除此段。}

  \textcolor{blue}{结论正文样式与文章正文相同：宋体、小四；行距：22 磅；间距段前段后均为 0 行。阅后删除此段。}
\end{conclusion}

% 参考文献：
% 添加文献时，请按 BibTeX 格式添加至 misc/ref.bib，并在正文所需位置使用 \cite{…} 引用。
% 如无特殊需要，无需改动相应 TeX 文件。
%%
% The BIThesis Template for Bachelor Graduation Thesis
%
% 北京理工大学毕业设计（论文）参考文献 —— 使用 XeLaTeX 编译
%
% Copyright 2020-2023 BITNP
%
% This work may be distributed and/or modified under the
% conditions of the LaTeX Project Public License, either version 1.3
% of this license or (at your option) any later version.
% The latest version of this license is in
%   http://www.latex-project.org/lppl.txt
% and version 1.3 or later is part of all distributions of LaTeX
% version 2005/12/01 or later.
%
% This work has the LPPL maintenance status `maintained'.
%
% The Current Maintainer of this work is Feng Kaiyu.
%
% Compile with: xelatex -> biber -> xelatex -> xelatex
%
% 如无特殊需要，本页面无需更改

\begin{bibprint}

\printbibliography[heading=none]

\end{bibprint}

% 附录：在相应 TeX 文件撰写；不需要时可删除
%%
% The BIThesis Template for Bachelor Graduation Thesis
%
% 北京理工大学毕业设计（论文）附录 —— 使用 XeLaTeX 编译
%
% Copyright 2020-2023 BITNP
%
% This work may be distributed and/or modified under the
% conditions of the LaTeX Project Public License, either version 1.3
% of this license or (at your option) any later version.
% The latest version of this license is in
%   http://www.latex-project.org/lppl.txt
% and version 1.3 or later is part of all distributions of LaTeX
% version 2005/12/01 or later.
%
% This work has the LPPL maintenance status `maintained'.
%
% The Current Maintainer of this work is Feng Kaiyu.
%
% Compile with: xelatex -> biber -> xelatex -> xelatex

\begin{appendices}
  附录相关内容…

  % 这里示范一下添加多个附录的方法：
  % 使用 \section 来添加一个附录

  \section{\LaTeX 环境的安装}
  \LaTeX 环境的安装。

  \section{BIThesis 使用说明}
  BIThesis 使用说明。

  \textcolor{blue}{附录是毕业设计（论文）主体的补充项目，为了体现整篇文章的完整性，写入正文又可能有损于论文的条理性、逻辑性和精炼性，这些材料可以写入附录段，但对于每一篇文章并不是必须的。附录依次用大写正体英文字母 A、B、C……编序号，如附录 A、附录 B。阅后删除此段。}

  \textcolor{blue}{附录正文样式与文章正文相同：宋体、小四；行距：22 磅；间距段前段后均为 0 行。阅后删除此段。}

\end{appendices}

% 致谢：在相应 TeX 文件撰写
%%
% The BIThesis Template for Graduate Thesis
%
% Copyright 2020-2023 Yang Yating, BITNP
%
% This work may be distributed and/or modified under the
% conditions of the LaTeX Project Public License, either version 1.3
% of this license or (at your option) any later version.
% The latest version of this license is in
%   http://www.latex-project.org/lppl.txt
% and version 1.3 or later is part of all distributions of LaTeX
% version 2005/12/01 or later.
%
% This work has the LPPL maintenance status `maintained'.
%
% The Current Maintainer of this work is Feng Kaiyu.

\begin{acknowledgements}

本论文的工作是在导师……。

\textcolor{blue}{
  致谢是对下列方面致谢：资助和支持者；协助完成研究工作和提供便利条件者；在研究工作中提出建议和提供帮助者；给予转载和引用权的资料、图片、文献、研究思想和设想的所有者；其他应感谢者。致谢语言要诚恳、恰当、简短。
}

\end{acknowledgements}


\end{document}
